% ===== wrap-up =====
\begin{frame}{이번 주차 정리}
  \begin{enumerate}
    \item \kb{합성곱 = 재사용의 연산} --- $n_{\text{out}}=\lfloor
      (n+2p-k)/s\rfloor+1$, params $(k^2C_{\text{in}}+1)C_{\text{out}}$,
      FLOPs에 $n_{\text{out}}^{\,2}$ 이 곱해진다. 미니 LeNet이
      MLP 대비 \textbf{0.12\%의 파라미터, 1/4의 연산}.
    \item \kb{좁아지고 깊어진다} --- 합성곱이 ``무엇을''(패턴),
      풀링이 ``얼마나 넓은 시야로''(수용장)를 결정. 깊이 열화는
      \textbf{최적화 문제}이고, 스킵 연결의 ``$+1$''이 그래디언트
      바닥을 깔아 \textbf{100층이 넘는 CNN}을 가능하게 함.
    \item \kb{같은 줄기, 다른 헤드} --- 분류(벡터 1개) / 탐지
      (박스 + NMS) / 분할(위치별 공유 1$\times$1 합성곱) /
      전이학습(앞쪽 동결, 헤드만 재학습). ``\textbf{앞쪽 = 범용,
      뒤쪽 = 특화}'' 비대칭이 전이학습의 핵심.
  \end{enumerate}
  \vskip0.8em
  \begin{block}{다음 주 --- Week 11. 시퀀스 모델}
    CNN이 가중치를 \textbf{공간}(이미지 위 위치)에 걸쳐
    재사용한다면, RNN은 같은 가중치를 \textbf{시간}(시퀀스의
    시점)에 걸쳐 재사용한다. 공간과 시간, 두 축에서
    ``\textbf{작은 단위 패턴을 재사용해 큰 구조를 쌓는다}''는
    동일한 원리를 본 셈이다 --- 이번 주차의 ``재사용''이
    바로 그 연결고리.
  \end{block}
\end{frame}
